\section{Reference Policies}
\label{sec:method}

The name \ksa{} simply abbreviates the three actions
(\keep/\supersede/\abstain); it is not a claim that the action vocabulary is
novel. The released pilot includes two policies over this vocabulary and two
shortcut probes. All commit actions pass through the same deterministic state
transition and invariant validator.

\subsection{Transparent reference policy}
\label{sec:policy}

The reference policy parses the raw perturbation text and prior state with
auditable lexical predicates. In priority order it (1) \abstain{}s on
unattributed third-party, laundered, garbled, or unconfirmed evidence; (2)
\abstain{}s on a hedged temporary candidate; (3) \keep{}s the durable value and
optionally stores a separate transient note for a clear temporary override; (4)
\supersede{}s when a usable durable candidate conflicts with the active value;
and (5) otherwise \abstain{}s. The complete predicates and value parser are
released in \texttt{benchmarks/mgp6/systems/}.

This policy is deterministic and inspectable, but it was written against the
same annotation manual and phrase generator as the pilot. It is therefore a
\emph{reference implementation of the policy}, not an independently learned
method. A five-feature weighted scorecard exists in the broader prototype code,
but it did not generate the reported MGP-6 predictions; we do not describe or
claim it as the evaluated method. This distinction removes an earlier
implementation--manuscript mismatch.

\subsection{State transition and lineage}

On \supersede{} the transition marks the prior record superseded, appends one
active successor, and points its \texttt{supersedes\_id} to the retained
predecessor. On \keep{} a clear temporary override may be a separate transient
record; on \abstain{} the modeled state is unchanged. The evaluator scores these
state properties independently of the action, so a correct action with corrupted
lineage is not counted as a correct state transition.

\subsection{Structured and unstructured LLM policies}

We also run the governance decision through a structured LLM policy prompted
with the typed state, candidate, evidence fields, and the three action
definitions, and an unstructured zero-shot policy that receives the raw event
and prior state without that scaffold. Both use
\texttt{qwen/qwen3.5-9b}. Because the prompt, candidate representation, and
pilot were author-designed together, even this comparison remains diagnostic.

\subsection{Shortcut probes}

The type-only and conflict-type-only probes learn a majority action for each
metadata value on train+dev and apply the lookup unchanged to test. They never
read the dialogue, prior value, or candidate evidence. Their purpose is to test
whether the collection rewards scenario recognition without memory reasoning.
